\documentclass[11pt]{article}

\usepackage[a4paper,margin=25mm,headheight=15pt]{geometry}
\usepackage[T1]{fontenc}
\usepackage[utf8]{inputenc}
\usepackage{lmodern}
\usepackage{microtype}
\usepackage{amsmath}
\usepackage{booktabs}
\usepackage{array}
\usepackage{xcolor}
\usepackage{fancyhdr}
\usepackage{tikz}
\usetikzlibrary{positioning,shapes.geometric}
\usepackage[hidelinks]{hyperref}

\definecolor{Accent}{HTML}{315B7D}
\definecolor{LightBlue}{HTML}{EAF2F8}
\definecolor{MidBlue}{HTML}{8EB4CE}
\definecolor{Muted}{HTML}{5B6570}
\definecolor{Aqu}{HTML}{3A7D77}
\definecolor{Pol}{HTML}{D08A45}
\definecolor{Grey}{HTML}{C7CDD2}

\hypersetup{
  pdftitle={Draft main-text Results and Discussion with figure plan},
  pdfauthor={Poikela et al.}
}

\pagestyle{fancy}
\fancyhf{}
\lhead{\small\color{Muted}Main-text planning draft}
\rhead{\small\color{Muted}Poikela et al.}
\cfoot{\small\thepage}
\renewcommand{\headrulewidth}{0.4pt}
\setlength{\parindent}{0pt}
\setlength{\parskip}{0.7em}
\renewcommand{\arraystretch}{1.18}

\newcommand{\Faq}{\textit{F.~aquilonia}}
\newcommand{\Fpo}{\textit{F.~polyctena}}

\title{\vspace{-3em}
\textbf{\color{Accent}Draft main-text Results and Discussion}\\[0.3em]
\Large Predictability and genomic organisation of ancestry sorting}
\author{Poikela et al.}
\date{\small Working draft excluding the epistasis screen and simulations}

\begin{document}
\maketitle
\vspace{-1.5em}

\begin{center}
\small\color{Muted}
The first two sections are written as candidate manuscript text. The final
section is an editorial figure and supplementary-material plan. Climate
associations are interpreted as conditional and exploratory.
\end{center}

\section*{Results}

\begin{figure}[t]
\centering
\includegraphics[width=\textwidth]{figures/main_ancestry_sorting.pdf}
\caption{\textbf{Predictability and genomic organisation of ancestry sorting.}
(\textbf{A}) Sort-class proportions among oriented parent-polymorphic loci:
near-fixation is predominantly unidirectional, with bidirectional fixation almost
absent. (\textbf{B}) Among unidirectionally sorted loci, the fraction fixing toward
\Faq{} rises across diagnostic-index deciles, crossing parity near the centre.
(\textbf{C}) Fraction sorted versus recombination at the SNP and LD-reduced-unit
levels; the unit-level collapse in low recombination reflects removal of spatial
pseudoreplication. (\textbf{D}) Diagnostic-index enrichment of climate-associated
clusters (percentage points per +10 percentage-point climate rate) for PC1 and PC2,
with model-based and chromosome block-bootstrap 95\% intervals; the nominal
enrichment is not robust to genomic dependence.}
\label{fig:main}
\end{figure}

\subsection*{Ancestry sorting is predominantly directional across hybrid populations}

To determine whether hybrid genome evolution repeatedly favoured the same
parental ancestry, we quantified near-fixation across twenty geographically
separated \Faq{} $\times$ \Fpo{} populations. The analysis extends the
three-population comparison of Nouhaud et al. (2022) in two respects: it
separates repeated fixation in the same parental direction from fixation in
opposing directions, and it accounts for local linkage rather than assigning
equal weight to fixed physical windows.

Among SNPs that remained polymorphic in the pooled parental samples, 67.96\%
were not near-fixed in at least half of the hybrid populations. Of those that
did show repeated sorting, 17.50\% sorted toward \Fpo{} and 14.40\% toward
\Faq{}, whereas only 0.14\% showed strong sorting in opposing directions among
populations (Fig.~\ref{fig:main}A). Thus, the main feature of the data was not
near-fixation itself, but the consistency of its parental direction across
populations. Because the populations may share aspects of their demographic
history, these proportions describe repeatability across the sampled panel
rather than proving independent evolutionary replication.

Marker-level summaries can overstate the genomic extent of repeated sorting
when many correlated markers report the same event. We therefore reduced
1,114,340 markers to 474,014 LD clusters and represented each cluster once.
Among clusters selected because they contained at least one SNP-level sorting
signal, aggregation removed the sorting classification increasingly often as
cluster size increased: 74.7\% of clusters containing at least five markers
were unsorted after aggregation, compared with 5.0\% of the smallest units.
This comparison is conditional on the presence of a marker-level signal and
does not estimate the genome-wide fraction of unsorted clusters. The loss of
the aggregate call was unlikely to result primarily from poor consensus
construction, because consensus fidelity remained high (median 0.856) and only
7 of 16,597 multi-marker unsorted clusters fell below the predefined fidelity
threshold. Repeated sorting was therefore distributed mainly across numerous
small genomic units rather than a limited set of large LD regions.

\subsection*{The parental direction of sorting changes along the diagnostic-index axis}

The modest genome-wide excess of sorting toward \Fpo{} concealed a strong
reversal in direction along the diagnostic-index axis. Among
unidirectionally sorted loci, the fraction sorting toward \Faq{} increased
monotonically from approximately 0.15 in the least diagnostic decile to 0.74
in the most diagnostic decile, crossing parity near the centre of the DI
distribution (Spearman \(\rho=0.31\); Fig.~\ref{fig:main}B). In a logistic model
including pooled parental MAF, DI strongly increased the probability of
\Faq-directed sorting (DI coefficient \(=0.082\), \(z=181\)), whereas higher
parental MAF weakly shifted sorting toward \Fpo{}.

Larger LD clusters were both more diagnostic and more likely to sort toward
\Faq{} in unadjusted summaries. This relationship did not reflect an
independent effect of cluster size. After DI and parental MAF were included,
the size coefficient reversed direction
\(\{\log_2(\mathrm{size})=-0.20,\ z=-22\}\), whereas the DI coefficient
remained nearly unchanged (\(0.083,\ z=120\)). Parental diagnosticity, rather
than the physical extent of a correlated region, was therefore the principal
axis organising the direction of repeated sorting.

\subsection*{Linkage changes the apparent abundance, but not the main direction, of sorting}

DI was nearly orthogonal to recombination rate
(\(\rho=-0.03\)) and pooled parental MAF (\(\rho=-0.02\)). At the analysed
SNPs, high DI was associated most strongly with lower parental expected
heterozygosity (\(\rho=-0.46\)) and only weakly with the between-parent
allelic-difference summary (\(\rho=0.13\)). The diagnostic-index gradient
therefore cannot be reduced to a low-recombination gradient or to elevated
between-parent difference alone.

The importance of representing LD-correlated markers once was clearest in
regions of lowest recombination. In the lowest recombination decile, 17\% of SNPs
were classified as sorted, but only 6\% of LD-reduced units were sorted,
compared with approximately 42\% of units across the remaining deciles
(Fig.~\ref{fig:main}C). The excess marker-level signal in low recombination
therefore largely reflected spatial pseudoreplication. Consistent with this
interpretation, aggregate sorting increased rather than decreased with
recombination rate.

Among unidirectionally sorted units, DI remained by far the strongest predictor
of parental direction in a model including recombination, parental MAF, and
cluster size (standardised DI coefficient \(=1.46,\ z=125\)). Recombination
made a much smaller contribution
(\(-0.09,\ z=-7.6\)), corresponding to a weak tendency toward \Faq{} ancestry
at lower recombination rates. This residual association is descriptive and is
not sufficient on its own to identify its evolutionary cause. Nevertheless,
both relationships were robust to residual genomic dependence. A chromosome
block bootstrap gave 95\% intervals of 1.41--1.52 for the DI coefficient and
\(-0.145\) to \(-0.037\) for the recombination coefficient; intervals from
10-cM genetic-map blocks were similar. The positive relationship between
recombination and sorting magnitude was likewise retained
(coefficient \(=0.052\), chromosome-block 95\% interval 0.046--0.057).

\subsection*{Climate association does not robustly explain diagnosticity or repeated directional sorting}

The hybrid populations span marked climatic variation. The first two
principal components of nineteen bioclimatic variables jointly explained 86\%
of this variation. PC1 primarily captured a temperature and seasonality
gradient, with the high-altitude Swiss population resembling colder northern
populations, whereas PC2 primarily separated populations by precipitation.
Genotype--environment associations were evaluated while accounting for
among-population covariance; because PC2 covaried with genome-wide ancestry and
this relationship was strongly influenced by \AA{}land, the
\AA{}land-excluded analysis with a fixed covariance matrix was treated as
primary.

Under model-based uncertainty, climate-associated LD clusters appeared to
contain a higher proportion of ancestry-diagnostic markers. Among clusters
with at least fifty markers, a 10-percentage-point increase in the proportion
of climate-associated members corresponded to increases of 4.5 percentage
points for PC1 (\(P=8\times10^{-4}\)) and 3.7 percentage points for PC2
(\(P=0.013\)) in the proportion of diagnostic markers. This nominal enrichment
was not robust to genomic dependence: chromosome block-bootstrap 95\%
intervals were \(-4.0\) to 13.1 for PC1 and \(-0.7\) to 17.8 for PC2, with
similar conclusions from 10-cM blocks (Fig.~\ref{fig:main}D).

The estimates were also attenuated after accounting for variables related to
allele-frequency variation and power. Adding mean hybrid MAF reduced the PC1
and PC2 coefficients to 3.5 and 2.5 percentage points, respectively. After
further adjustment for mean recombination rate and BayPass XtX, a measure of
among-population differentiation, the corresponding coefficients were 1.8 and
1.7 percentage points (\(P=0.16\) and 0.23); chromosome block-bootstrap
intervals again included zero. Because XtX and the climate Bayes factor are
derived from the same allele-frequency data, this is a conservative
sensitivity analysis rather than an estimate of an independently adjusted
climate effect. Together, the analyses show that the nominal enrichment is
compatible with shared dependence on allele-frequency differentiation and
cannot be identified as a climate-specific association.

The same climate-association rate did not predict repeated directional
sorting. In clusters with at least fifty markers, the odds ratios were 0.99
for PC1 and 1.03 for PC2, with \(P>0.95\) for both. This absence of enrichment
should be interpreted cautiously: genotype--environment association requires
allele-frequency variation among populations, whereas repeated sorting toward
the same parent removes that variation. Nevertheless, the contrast is
informative: the measured climatic gradients did not robustly explain either
diagnostic content beyond general differentiation or the repeated direction
of ancestry sorting.

\section*{Discussion}

Across twenty hybrid populations, ancestry sorting was strikingly asymmetric:
loci that repeatedly approached fixation almost always did so toward the same
parent, while strong fixation in opposing directions was rare. This extends
the repeatable sorting described across three populations by Nouhaud et al.
(2022) to a much broader population panel. The additional populations also
allow repeatability to be decomposed explicitly into consistent and opposing
directions, although we remain cautious about treating all populations as
fully independent evolutionary replicates.

The central result is not a uniform bias toward one parent. Instead, the
parental direction of sorting reverses along the diagnostic-index axis:
low-DI loci tend to sort toward \Fpo{}, whereas high-DI loci tend to sort
toward \Faq{}. This explains why analyses restricted to strongly diagnostic
loci can yield an apparent \Faq{} excess while the broader marker set shows a
slight \Fpo{} excess. DI is not simply a proxy for recombination or cluster
size. At the analysed SNPs it is associated mainly with reduced parental
expected heterozygosity, suggesting that the direction of sorting is organised
by the parental frequency landscape rather than by a single class of large,
low-recombination regions.

Accounting for LD changes the biological picture. Marker-level analyses make
sorting appear concentrated in low-recombination regions because many nearby
markers repeat the same ancestry signal. Once each LD cluster contributes once,
sorting is least frequent in the lowest-recombination decile and is carried
primarily by numerous small units. The weak residual association between low
recombination and \Faq-directed sorting is therefore secondary to the much
stronger diagnostic-index gradient. More generally, this comparison shows that
the apparent genomic extent of repeatable evolution can depend strongly on
whether markers or correlated genomic units are counted.

Climate association did not provide a robust explanation for repeated
sorting. Directionally sorted units were not enriched in climate-associated
regions, and the nominal enrichment of ancestry-diagnostic markers disappeared
once residual genomic dependence was respected. Its attenuation after
accounting for hybrid MAF and XtX further indicates that it could arise from
the allele-frequency differentiation that contributes to both diagnosticity
and BayPass evidence. XtX is a deliberately strong control and may remove
biologically relevant as well as confounding variation, so this sensitivity
analysis does not demonstrate that climate is irrelevant. Moreover, locally
varying selection could drive populations in different directions, whereas
repeated near-fixation reduces the among-population variation used to detect
genotype--environment association. The present data nevertheless provide no
robust evidence that the measured climatic gradients organise the repeated
direction of ancestry sorting.

Together, the results identify diagnosticity as the main organising axis of
early hybrid genome evolution in this system. Recombination determines how
many markers repeat a signal and makes smaller but block-robust contributions
to sorting magnitude and direction. In contrast, neither diagnostic content
nor repeated directional near-fixation showed a robust climate-specific
association. The emerging picture is therefore one of many small ancestry
units sorting in a highly structured parental direction, rather than a few
large genomic regions responding uniformly to linkage or the measured
climatic gradients. These conclusions remain associative: shared population
history and conditional testing limit causal interpretation, and individual
candidate regions require independent biological scrutiny.

\section*{Editorial figure plan}

\subsection*{Recommended main figure: one integrated ancestry-sorting figure}

The main figure should carry the central narrative rather than the validation
machinery. A four-panel figure is sufficient:

\noindent This layout is realised in Figure~\ref{fig:main} (shown in the Results),
which progresses from the phenomenon (A), to its principal predictor (B), to the
effect of genomic linkage on its apparent abundance (C), and finally to its
relationship with environmental association (D).

Panel A establishes the phenomenon. Panel B is the biological centre of the
figure and should receive the most visual space. Panel C demonstrates why the
LD reduction changes the interpretation. Panel D closes the extrinsic-climate
branch without requiring separate Manhattan plots in the main text.

\subsection*{Material better placed in supplementary figures and tables}

\begin{tabular}{>{\raggedright\arraybackslash}p{0.28\textwidth}
                >{\raggedright\arraybackslash}p{0.62\textwidth}}
\toprule
Supplementary item & Recommended content \\
\midrule
LD decay diagnostics &
Chromosome-specific decay fits, ROC evaluation of local LD support, and
example genomic tracks. \\
\addlinespace
Complexity-reduction validation &
Stage 1 versus combined clustering in extended-LD regions, fidelity
distributions, threshold sensitivity, residual inter-unit LD, and a concise
parameter table. \\
\addlinespace
Sorting robustness &
Full class counts, conditional aggregation by cluster size, fidelity check for
apparently diluted clusters, exact-test summaries, and sensitivity to sorting
thresholds. \\
\addlinespace
Genomic architecture &
Full parental-summary and recombination decile plots, complete regression
tables, and consensus-versus-representative validation. \\
\addlinespace
Climate association &
Ancestry--climate confounding check, nested cluster-size strata, the MAF/XtX
sensitivity ladder, block-aware uncertainty, and full model tables. \\
\addlinespace
Dependence-aware uncertainty &
Chromosome- and 10-cM-block confidence intervals for the principal sorting,
architecture, and climate coefficients, with the reconstruction checks and
bootstrap diagnostics. \\
\bottomrule
\end{tabular}

\subsection*{Suggested order in the full paper}

The co-authors' sampling, ancestry, and climate overview should establish the
system first. The text drafted here can then enter as one continuous section:
(1) directional repeatability, (2) diagnostic-index polarity, (3) genomic
architecture after LD reduction, and (4) overlap with climate association.
This order lets the main figure function as a bridge from the descriptive
population overview to the later mechanistic analyses without exposing the
internal analysis-module structure.

\section*{Reference}

Nouhaud P, Martin SH, Portinha B, Sousa VC, Kulmuni J. Rapid and predictable
genome evolution across three hybrid ant populations. \textit{PLoS Biology}.
2022;20:e3001914. \url{https://doi.org/10.1371/journal.pbio.3001914}.

\end{document}
